\documentclass[11pt]{article}

\usepackage[utf8]{inputenc}
\usepackage[T1]{fontenc}
\usepackage[spanish,es-nodecimaldot]{babel}
\usepackage{mathptmx}
\usepackage[margin=2.3cm]{geometry}
\usepackage{microtype}
\usepackage{amsmath}
\usepackage{enumitem}
\usepackage{titlesec}
\usepackage{parskip}

\titleformat{\section}{\normalfont\bfseries\large}{}{0pt}{}
\titlespacing*{\section}{0pt}{12pt}{4pt}
\titleformat{\subsection}{\normalfont\bfseries}{}{0pt}{}
\titlespacing*{\subsection}{0pt}{8pt}{2pt}
\setlist[itemize]{leftmargin=1.2em,itemsep=2pt,topsep=2pt}
\setlist[enumerate]{leftmargin=1.6em,itemsep=3pt,topsep=3pt}
\pagestyle{empty}
\setlength{\parskip}{5pt}
\hyphenpenalty=10000
\exhyphenpenalty=10000
\emergencystretch=2em

\begin{document}

\begin{center}
	{\small SMAII, Redes Complejas Multicapa}\\[6pt]
	{\Large\bfseries Proyecto final: sección de Antecedentes}\\[3pt]
	{\normalsize Efecto de una campaña sobre la opinión colectiva y predicción de su conveniencia}\\[4pt]
	{\small Septiembre de 2026}
\end{center}

\section*{Integrantes del equipo}
\begin{itemize}
	\item Jorge Miyoshi Mikami Alonso
	\item Luis Manuel Reyes Colín
\end{itemize}

\section*{Antecedentes}

Los primeros modelos que tratan la opinión de un grupo como algo que evoluciona en el tiempo
aparecen a mediados del siglo XX, y desde los ochenta la física estadística se sumó al tema bajo
el nombre de sociofísica. El punto de partida de nuestro proyecto es más reciente: en el año 2000
Deffuant y Weisbuch propusieron el modelo de confianza acotada \cite{Deffuant2000}, donde cada
persona tiene una opinión entre 0 y 1 y dos personas solo se influyen si sus opiniones están más
cerca que un umbral $\epsilon$. Con $\epsilon$ chica la población se fragmenta en grupos que ya no
se hablan; con $\epsilon$ grande todos terminan de acuerdo.

Pansanella \emph{et al.} \cite{Pansanella2023} le agregan a ese modelo el sesgo algorítmico (un
parámetro $\gamma$ que hace más probable que te toque hablar con quien ya piensa como tú) y medios
masivos, que meten como agentes ``tercos'' con opinión fija a los que cada agente se asoma con
probabilidad $p_m$. Prueban cuatro escenarios (un medio moderado, uno extremista, dos polarizados
y uno balanceado) y sus resultados son contraintuitivos: una población abierta es \emph{más} fácil
de manipular, un mensaje moderado repetido demasiado favorece minorías extremistas, y el sesgo
algorítmico acaba protegiendo a la población de radicalizarse por completo.

Yin \emph{et al.} \cite{Yin2025} atacan lo que el trabajo anterior no toca: la conducta. Arman una
red de dos capas, una de opinión sobre la economía con un modelo de umbral y otra de consumo tipo
SIS (Susceptible--Infectado--Susceptible, el modelo epidémico donde la conducta se contagia entre
vecinos pero no deja inmunidad, así que cada persona puede activarse y volver a dejar de consumir
cuantas veces sea), acopladas en ambos sentidos: más consumo baja el umbral para volverse
optimista, y más optimismo activa más consumo. Esa retroalimentación hace que el sistema salte de golpe entre un
estado pesimista de bajo consumo y uno optimista de alto consumo.

Lo que notamos pendiente es que Pansanella \emph{et al.} trabajan en campo medio, sin estructura de
red y sin ninguna consecuencia conductual de la opinión, mientras que Yin \emph{et al.} usan
opiniones binarias, sin confianza acotada ni mensaje externo. Sobre todo, ninguno de los dos
convierte su barrido de parámetros en algo que sirva para decidir de antemano. Nuestro aporte es
juntar ambos modelos en un esquema de dos capas y usar aprendizaje automático para entrenar un
predictor que, a partir de las condiciones iniciales y las características del mensaje, anticipe
si la campaña será favorable o contraproducente.

\section*{Objetivos de investigación}

\subsection*{Objetivo general}

Investigar en qué condiciones una campaña masiva logra mover la opinión de una población hacia la
postura que promueve y en cuáles resulta inútil o contraproducente, usando un modelo de confianza
acotada con sesgo algorítmico y medios masivos acoplado a una segunda capa de conducta, con el fin
de construir un criterio que permita anticipar si conviene lanzarla.

\subsection*{Objetivos particulares}

\begin{enumerate}
	\item Hacer la búsqueda bibliográfica y redactar los antecedentes del proyecto.

	\item Programar el modelo de confianza acotada con sesgo algorítmico y medios, y validarlo
	      replicando los resultados de Pansanella \emph{et al.}

	\item Construir las redes de simulación y agregar la segunda capa de conducta, definiendo
	      cómo se retroalimentan entre sí.

	\item Recorrer el espacio de parámetros ($\epsilon$, $\gamma$, $p_m$, $x_m$) y calcular las
	      métricas del estado final: número de grupos de opinión, tamaño del grupo alrededor de
	      la campaña y nivel de polarización.

	\item Clasificar los estados finales y entrenar un modelo que prediga si la campaña resulta
	      favorable o contraproducente.

	\item Analizar los resultados y comparar las conclusiones con los artículos de referencia.
\end{enumerate}

\begin{thebibliography}{9}
\setlength{\itemsep}{3pt}

\bibitem{Deffuant2000}
G. Deffuant, D. Neau, F. Amblard and G. Weisbuch, \emph{Mixing beliefs among interacting agents}, Advances in Complex Systems \textbf{3} (2000) 87--98.

\bibitem{Pansanella2023}
V. Pansanella, A. Sîrbu, J. Kertesz and G. Rossetti, \emph{Mass media impact on opinion evolution in biased digital environments: a bounded confidence model}, Scientific Reports \textbf{13} (2023) 14600.

\bibitem{Yin2025}
L. Yin, J. Chen, F. Gao and X. Wu, \emph{Coevolution of opinion and consumption behavior under a two-layer network framework}, Physica A \textbf{680} (2025) 131020.

\end{thebibliography}

\end{document}
